% Auto-typeset from the previously authored formula-by-formula study manuscript.
% Based on the thesis main branch f80e5ac903dd, 2026-10-31.

\chapter{矩阵 RHP、留数、无反射有限维化与三场重构}

对照正文：../hmb-nzbc-thesis/section/04-rhp.tex  \par
{\sloppy\raggedright 核心标签：thm:RHP，eq:Mdef，eq:jump，eq:resplus，eq:resminus，eq:cpair，prop:finite，eq:finite-schur，prop:theta，prop:reconstruction。\par}

本章真正的核心问题：\textbf{已经知道哪些 Jost 列解析，怎样把它们拼成一个可求解的矩阵函数 M？} 若反射为零，为什么这种复分析问题最终只剩下有限维线性方程组？\par

\section{为什么不直接取 \texorpdfstring{$M=\mu_-$}{M=mu-} 或 \texorpdfstring{$M=\mu_+$}{M=mu+}？}

上一章只证明\par

\[
\mu_{-,1},\ \mu_{+,2}\text{ 在 }D_+\text{ 解析};\qquad
\mu_{+,1},\ \mu_{-,2}\text{ 在 }D_-\text{ 解析}.
\]

因此，如果希望 M 在 $D_{+}$ 作为整个矩阵解析，第一列必须从 $\mu{}_{-}$ 第一列选，第二列必须从 $\mu{}_{+}$ 第二列选；$D_{-}$ 则反过来。\par

另一方面，透射系数 $s_{11}$、$s_{22}$ 的零点会产生离散谱。在 RHP 中通常将透射系数作为适当的归一化因子；这样一来，离散零点处自然出现简单极点，正好能够作为留数条件处理。\par

所以 eq:Mdef 为\par

\[
M(w)=
\begin{cases}
M_+(w)=\left(\dfrac{\mu_{-,1}}{s_{11}},\ \mu_{+,2}\right),
&w\in D_+,\\[6pt]
M_-(w)=\left(\mu_{+,1},\ \dfrac{\mu_{-,2}}{s_{22}}\right),
&w\in D_-.
\end{cases}
\]

这里括号表示“把两个列向量并排组成一个 $2\times{}2$ 矩阵”，不是普通有序对。$D_{+}$、$D_{-}$ 与轮廓方向由第二章确定，+、- 不是固定指上半平面和下半平面。\par

为何要除以 $s_{11}$、$s_{22}$？这一选择既使 M 在无穷远具有统一归一化，也让两个域的行列式恰好都等于背景 $\gamma{}$；下面立即验证。\par

\section{自己从两条散射列关系推出来}
\noindent{\small\textbf{对应原文：}\texttt{eq:jump}}\par

上一章的同时归一化散射关系 $\Psi{}_{-}=\Psi{}_{+}S$ 经过去相位后是\par

\[
\begin{aligned}
\mu_{-,1}
&=s_{11}\mu_{+,1}
+s_{21}e^{2\mathrm{i}\theta}\mu_{+,2},\\
\mu_{-,2}
&=s_{12}e^{-2\mathrm{i}\theta}\mu_{+,1}
+s_{22}\mu_{+,2}.
\end{aligned}
\]

先处理第二式。两边除以 $s_{22}$：\par

\[
\dfrac{\mu_{-,2}}{s_{22}}
=\dfrac{s_{12}}{s_{22}}e^{-2\mathrm{i}\theta}\mu_{+,1}
+\mu_{+,2}.
\]

由反射定义 $\widetilde r=s_{12}/s_{22}$，和 $M_{-}$ 第一列、第二列的定义，\par

\[
M_{-,2}
=\widetilde r e^{-2\mathrm{i}\theta}M_{-,1}
+M_{+,2}.
\]

所以\par

\[
M_{+,2}=M_{-,2}-\widetilde r e^{-2\mathrm{i}\theta}M_{-,1}.
\]

再处理第一式，除以 $s_{11}$：\par

\[
M_{+,1}
=\mu_{+,1}+r e^{2\mathrm{i}\theta}\mu_{+,2}.
\]

代入上一步的 $\mu{}_{+},_{2}=M_{+},_{2}$：\par

\[
\begin{aligned}
M_{+,1}
&=M_{-,1}
+r e^{2\mathrm{i}\theta}
\left(M_{-,2}-\widetilde r e^{-2\mathrm{i}\theta}M_{-,1}\right)\\
&=(1-r\widetilde r)M_{-,1}
+r e^{2\mathrm{i}\theta}M_{-,2}.
\end{aligned}
\]

两列合并：\par

\[
M_+=M_-J,
\qquad
J=\begin{pmatrix}
1-r\widetilde r&-\widetilde r e^{-2\mathrm{i}\theta}\\
r e^{2\mathrm{i}\theta}&1
\end{pmatrix}.
\]

检查 det J：\par

\[
\det J
=(1-r\widetilde r)
-(-\widetilde r e^{-2\mathrm{i}\theta})(r e^{2\mathrm{i}\theta})
=1.
\]

这就是完整的跳跃矩阵来源；没有必要机械记忆其四个元素。\textbf{先计算两列，再读出每列在旧基下的系数}，是所有此类 $2\times{}2 RHP$ 的通用方法。\par

\section{为什么 \texorpdfstring{$\det M=\gamma$}{det M=gamma} 而不是 1？}
\noindent{\small\textbf{对应原文：}\texttt{eq:detM}}\par

对 $D_{+}$：\par

\[
\begin{aligned}
\det M_+
&=\det\left(\dfrac{\mu_{-,1}}{s_{11}},\mu_{+,2}\right)\\
&=\dfrac{\det(\mu_{-,1},\mu_{+,2})}{s_{11}}\\
&=\gamma(w),
\end{aligned}
\]

因为上一章 $s_{11}=det(\mu{}_{-},_{1},\mu{}_{+},_{2})/\gamma{}$。$D_{-}$ 也同样成立。因此\par

\[
\det M(w)=\gamma(w)=1+\dfrac{Q^2}{w^2}.
\]

特别重要：在 $w=\pm{}iQ$，$\gamma{}=0$，因此 M 在这些点可能\textbf{解析但不满秩}。不能像普通单位行列式 RHP 那样直接说 $M^{-1}$ 在所有非极点都解析。第五章需要仔细证明 Lax 对数导数在分支点的表观奇性可去。\par

\section{eq:Minfty 与 eq:Mzero：为什么有两个不同的无穷远？}

上一章高能展开给出\par

\[
M(w)=I+\dfrac{M^{(1)}}{w}+O(w^{-2}),
\qquad w\to\infty.
\]

由于非零背景的两叶已经单值化到 w 平面，原谱曲面的另一无穷远对应 $w=0$。换叶对称将 $w\to{}\infty{}$ 的行为转为\par

\[
M(w)=-\dfrac{\mathrm{i}Q}{w}\sigma_1+O(1),
\qquad w\to0.
\]

千万不要把第二个式子解释为“新出现一个未知的离散极点”。它是\textbf{由非零背景固定的主项}，无论有多少孤子都应保留。\par

这两个归一化也解释了无反射部分分式表示中为何一开始就有 $I-iQ\sigma{}_{1}/w$，而不是简单 I 加若干孤子极点。\par

\section{两种约化如何限制 M？}
\noindent{\small\textbf{对应原文：}\texttt{eq:Msym}}\par

上一章得到两种基本变换：\par

\begin{enumerate}[leftmargin=2em]
\item 共轭变换 w\ensuremath{\to}\textbackslash{}overline w，伴随 K 矩阵。
\item 换叶变换 w\ensuremath{\to}h(w)=\ensuremath{-}Q\textsuperscript{2}/w，伴随 C(w)=iw\ensuremath{\sigma}\textsubscript{1}/Q。
\end{enumerate}

于是\par

\[
M(w)=K\overline{M(\overline w)}K^{-1},
\qquad
M(h(w))=M(w)\dfrac{\mathrm{i}w}{Q}\sigma_1.
\]

为什么不只是对散射矩阵 S 有对称？因为 M 的每一列都是特定的 Jost 解析列（加透射因子），这些变换会按规则交换列以及 $D_{+}$、$D_{-}$。只要在两侧用相同的拼接定义，就能继承约化。\par

它们的作用：\par

\begin{itemize}[leftmargin=2em]
\item 限定一组允许的离散点，不是每个 \ensuremath{\zeta} 可以自由选。
\item 把一处留数的系数关联到另外三处。
\item 决定无反射矩阵 M 的构造是否与真实 HMB 谱问题相容。
\end{itemize}

\section{一个透射零点如何产生第一列极点？}
\noindent{\small\textbf{对应原文：}\texttt{eq:resplus}}\par

设 $\zeta{}\in{}D_{+}$，$s_{11}(\zeta{})=0$、$s'_{11}(\zeta{})\ne{}0$。Taylor 展开为\par

\[
s_{11}(w)=s_{11}'(\zeta)(w-\zeta)+O((w-\zeta)^2).
\]

M 第一列是 $\mu{}_{-},_{1}/s_{11}$，因此\par

\[
\operatorname*{Res}_{w=\zeta}M_1(w)
=\dfrac{\mu_{-,1}(\zeta)}{s_{11}'(\zeta)}.
\]

而在透射零点，$\Psi{}_{-}$ 第一列与 $\Psi{}_{+}$ 第二列线性相关：\par

\[
\Psi_{-,1}(\zeta)=b\,\Psi_{+,2}(\zeta).
\]

去振荡时不要忘记两列的指数方向相反：\par

\[
\mu_{-,1}(\zeta)
=b e^{2\mathrm{i}\theta(\zeta)}\mu_{+,2}(\zeta).
\]

把 $c=b/s'_{11}(\zeta{})$ 和 $C=ce^{2i\theta{}(\zeta{})}$ 代入：\par

\[
\boxed{\operatorname*{Res}_{w=\zeta}M_1(w)=C\,M_2(\zeta).}
\]

此点 $M_{2}$ 正则，所以其留数为 0。这说明极点只出现在矩阵的第一列，而系数由第二列的值决定。\par

\section{为什么共轭点有负号？}
\noindent{\small\textbf{对应原文：}\texttt{eq:resminus}}\par

在共轭点 $\overline \zeta{}$，透射零点位于 $D_{-}$，极点改在 M 的第二列。由 K 共轭约化，RHP 的对应留数为\par

\[
\operatorname*{Res}_{w=\overline\zeta}M_2(w)
=-\overline C\,M_1(\overline\zeta),
\qquad
\operatorname*{Res}_{w=\overline\zeta}M_1(w)=0.
\]

\textbf{负号来自 K 的具体矩阵结构}，不是任意约定。想检验，令某个局部主部为 $(v,0)/(w-\zeta{})$，先逐项共轭，再左右乘 K 和 $K^{-1}$，观察第一列与第二列如何交换。\par

把留数条件写成矩阵形式也很有用：\par

\[
\operatorname*{Res}_{w=\zeta}M(w)
=\lim\limits_{w\to\zeta}M(w)
\begin{pmatrix}0&0\\C&0\end{pmatrix}.
\]

右侧的极点矩阵为幂零矩阵 A，满足 $A^{2}=0$。虽然 M 本身有极点，但 MA 中只保留 M 的正则第二列，因此极限有意义。这种结构是第五章证明对数导数中离散极点可去的关键。\par

\section{换叶留数常数的因子怎样来的？}
\noindent{\small\textbf{对应原文：}\texttt{eq:cpair}}\par

取同一 $D_{+}$ 中的另一个配对点\par

\[
\widehat\zeta=h(\overline\zeta)
=-\dfrac{Q^2}{\overline\zeta}.
\]

令 $v=h(w)$。在 $w=\overline\zeta{}$ 附近，Taylor 展开：\par

\[
v-\widehat\zeta
=h'(\overline\zeta)(w-\overline\zeta)
+O((w-\overline\zeta)^2),
\quad
h'(\overline\zeta)=\dfrac{Q^2}{\overline\zeta^{\,2}}.
\]

若一个函数在 $w_{0}$ 有简单极点，其在新坐标 v 的留数要乘以 $h'(w_{0})$。还需结合换叶对称中的列交换及相位 $\theta{}(h(w))=-\theta{}(w)$。最终得到\par

\[
\boxed{
c_{\widehat\zeta}
=-\dfrac{Q^2}{\overline\zeta^{\,2}}\overline{c_\zeta}.
}
\]

请注意，不仅有共轭，还有导数因子 $Q^{2}/\overline\zeta{}^{2}$；直接把 $c_{\hat \zeta{}}$ 写成 $-\bar c_\zeta{}$ 是错误的。第五章选择 $c_{1}$ 后，$c_{2}$ 就被这一关系锁定。\par

\section{RHP 的局部条件为何不能略掉？}

第二章的谱曲线在 $w=\pm{}Q$ 交叉，$w=\pm{}iQ$ 为背景特征向量退化点，$w=w_{a}\pm{}$ 为共振原像。Jost 解只在普通点上有前文的解析结论。若研究一般含反射 RHP，就必须另外说明这几类点的局部行为。\par

因此原文定理 thm:RHP \textbf{主要由已知场的直接散射导出}。它并不等同于对任意给定的 r、$c_{j}$ 都有唯一解的逆散射定理。\par

第五章的一轨道无反射显式矩阵，会直接核查这些特殊点，没有理由把一般含反射情形的困难自动抹掉。\par

\section{无反射时为什么可以做部分分式？}
\noindent{\small\textbf{对应原文：}\texttt{prop:finite}}\par

当 $r=\widetilde r=0$，刚才得到的 J 变成单位矩阵。若轮廓非极点上的侧值能解析拼接，M 就成为除指定孤子极点和固定 $w=0$ 主项外的亚纯函数。\par

扣除所有极点主部，剩余矩阵整解析且无穷远趋于 I。按 Liouville 定理剩余就是 I，于是 eq:partialfraction：\par

\[
\boxed{
\begin{aligned}
M(w)=I-\dfrac{\mathrm{i}Q}{w}\sigma_1
+\sum\limits_{j=1}^{m}\left(
\dfrac{C_jv_je_1^{\mathsf T}}{w-\zeta_j}
-\dfrac{\overline{C_j}u_je_2^{\mathsf T}}{w-\overline{\zeta_j}}
\right).
\end{aligned}}
\]

这里\par

\[
v_j=M_2(\zeta_j),\qquad
u_j=M_1(\overline{\zeta_j}).
\]

注意：$u_j$、$v_j$ 是二维\textbf{向量}，不是场 u、v；$e_{1}^{T}=(1,0)$、$e_{2}^{T}=(0,1)$ 是行向量。乘上 $e_{1}^{T}$ 意味着这一项只进入第一列；乘上 $e_{2}^{T}$ 只进入第二列。\par

为什么正域和共轭点有不同符号？正域留数是 $+C_{j}v_{j}$，负域共轭留数是 $-\overline C_{j}u_{j}$，因此两种主部本来就带不同符号。\par

\section{在 \texorpdfstring{$\zeta_j$}{zetaj} 取第二列为何得到第一组方程？}
\noindent{\small\textbf{对应原文：}\texttt{eq:finite1}}\par

观察部分分式的第二列：\par

\begin{itemize}[leftmargin=2em]
\item 固定背景 I 的第二列为 e\textsubscript{2}；
\item \ensuremath{-}iQ\ensuremath{\sigma}\textsubscript{1}/w 的第二列为 \ensuremath{-}iQe\textsubscript{1}/w；
\item 第一组 e\textsubscript{1}\textsuperscript{T} 主部仅有第一列，\textbf{对第二列没有贡献}；
\item 第二组 e\textsubscript{2}\textsuperscript{T} 主部的第二列为 \ensuremath{-}\textbackslash{}overline C\_k u\_k/(w\ensuremath{-}\textbackslash{}overline\ensuremath{\zeta}\_k)。
\end{itemize}

在 $w=\zeta{}_j$ 取值（这个点第二列正则），得到\par

\[
\boxed{
v_j=e_2-\dfrac{\mathrm{i}Q}{\zeta_j}e_1
-\sum\limits_{k=1}^{m}
\dfrac{\overline{C_k}u_k}{\zeta_j-\overline{\zeta_k}}.
}
\]

这是 eq:finite1。分母不为零，因为 $\zeta{}_j$ 属于 $D_{+}$、$\overline\zeta{}_k$ 属于 $D_{-}$ 且不重合。\par

\section{在共轭点取第一列为何得到第二组方程？}
\noindent{\small\textbf{对应原文：}\texttt{eq:finite2}}\par

同样，部分分式第一列为\par

\[
M_1(w)=e_1-\dfrac{\mathrm{i}Q}{w}e_2
+\sum\limits_{k=1}^{m}\dfrac{C_kv_k}{w-\zeta_k}.
\]

在 $w=\overline\zeta{}_j$ 取值：\par

\[
\boxed{
u_j=e_1-\dfrac{\mathrm{i}Q}{\overline{\zeta_j}}e_2
+\sum\limits_{k=1}^{m}
\dfrac{C_kv_k}{\overline{\zeta_j}-\zeta_k}.
}
\]

这就是 eq:finite2。\textbf{到这里，未知的解析函数 M 已经转化为有限个未知二维向量 u\_j、v\_j}。接下来只需线性代数。\par

\section{Schur 补怎么得到？}
\noindent{\small\textbf{对应原文：}\texttt{eq:finite-schur}}\par

定义标量系数矩阵\par

\[
\mathcal A_{jk}
=\dfrac{\overline{C_k}}{\zeta_j-\overline{\zeta_k}},
\qquad
\mathcal B_{jk}
=\dfrac{C_k}{\overline{\zeta_j}-\zeta_k}.
\]

定义已知向量：\par

\[
f_j=e_2-\dfrac{\mathrm{i}Q}{\zeta_j}e_1,
\qquad
g_j=e_1-\dfrac{\mathrm{i}Q}{\overline{\zeta_j}}e_2.
\]

将各向量按 j 排列成块向量，方程成为\par

\[
\boldsymbol v=\boldsymbol f-\mathcal A\boldsymbol u,\qquad
\boldsymbol u=\boldsymbol g+\mathcal B\boldsymbol v.
\]

这里 $\mathcal A$、$\mathcal B$ 是 $m\times{}m$ 标量矩阵，每个标量乘的是一个二维向量。换言之，同一个 $m\times{}m$ 系数矩阵分别处理二维向量的第一分量和第二分量。\par

将第二式代入第一式：\par

\[
\begin{aligned}
\boldsymbol v
&=\boldsymbol f-\mathcal A(\boldsymbol g+\mathcal B\boldsymbol v)\\
&=\boldsymbol f-\mathcal A\boldsymbol g-\mathcal A\mathcal B\boldsymbol v.
\end{aligned}
\]

于是\par

\[
(I_m+\mathcal A\mathcal B)\boldsymbol v
=\boldsymbol f-\mathcal A\boldsymbol g.
\]

若\par

\[
\det(I_m+\mathcal A\mathcal B)\ne0,
\]

就有 eq:finite-schur：\par

\[
\boxed{
\boldsymbol v=(I_m+\mathcal A\mathcal B)^{-1}
(\boldsymbol f-\mathcal A\boldsymbol g),
\qquad
\boldsymbol u=\boldsymbol g+\mathcal B\boldsymbol v.
}
\]

\textbf{这是代数系统的可逆性条件，不等价于完整 RHP 的全部条件}。解出的矩阵还必须检查谱约化、行列式、阈值行为与端态相容性。第五章正是用一个特殊四点轨道把这些条件逐一补齐。\par

\section{为什么会有迹公式？}
\noindent{\small\textbf{对应原文：}\texttt{prop:theta}}\par

无反射时在普通轮廓上 $s_{11}s_{22}=1$。把 $D_{+}$ 的 $s_{11}$ 与 $D_{-}$ 的 $\dfrac{1}{s}_{22}$ 拼接，若没有其他奇异点，所得函数可由有限零点和极点表示。\par

除去这些零、极点后是整个黎曼球面上无零无极点且无穷远值为 1 的函数，故常等于 1。于是\par

\[
s_{11}(w)
=\prod\limits_{j=1}^{m}
\dfrac{w-\zeta_j}{w-\overline{\zeta_j}},
\qquad w\in D_+.
\]

这是条件性无反射迹公式 eq:trace。它本质上是“利用零点、极点与无穷远归一化确定有理函数”。如果有一般反射，迹公式通常还会包含轮廓上的对数积分，不能直接沿用这个纯乘积。\par

\section{同相端态怎样限制离散谱？}
\noindent{\small\textbf{对应原文：}\texttt{eq:theta-condition}}\par

对本文两端\textbf{完全相同}的背景，空间散射在 $w=0$ 的对应归一化为单位。由上面乘积在 $w=0$ 代入：\par

\[
s_{11}(0)
=\prod\limits_{j=1}^{m}
\dfrac{-\zeta_j}{-\overline{\zeta_j}}
=\prod\limits_{j=1}^{m}\dfrac{\zeta_j}{\overline{\zeta_j}}.
\]

所以同相背景要求\par

\[
\boxed{
\prod\limits_{j=1}^{m}\dfrac{\zeta_j}{\overline{\zeta_j}}=1.
}
\]

若一完整轨道在 $D_{+}$ 包含 $\zeta{}$ 与 $h(\overline\zeta{})$，将二者的比值相乘：\par

\[
\dfrac{\zeta}{\overline\zeta}\,
\dfrac{h(\overline\zeta)}{\overline{h(\overline\zeta)}}
=\left(\dfrac{\zeta}{\overline\zeta}\right)^2.
\]

因为 $h(\overline\zeta{})=-Q^{2}/\overline\zeta{}$，且 $\overline{h(\overline\zeta{})}=-Q^{2}/\zeta{}$。令 $\zeta{}=|\zeta|e^{i\varphi{}}$，则 $\zeta{}/\bar\zeta{}=e^{2i\varphi{}}$。于是 N 条轨道合起来是\par

\[
e^{4\mathrm{i}\sum\limits_{k=1}^{N}\arg\zeta_k}=1,
\]

即\par

\[
4\sum\limits_{k=1}^{N}\arg\zeta_k
\equiv0\pmod{2\pi}.
\]

纯虚外部上半平面代表的相角是 $\pi{}/2$，因此其每条轨道贡献 $e^{2\pi{}i}=1$，自动满足。这解释第五章为什么选纯虚谱点：它能自然保持两端同相，而一般复谱点需要额外约束。\par

\section{光场为何从 \texorpdfstring{$w\to\infty$}{w -> infinity} 的系数读出？}
\noindent{\small\textbf{对应原文：}\texttt{prop:reconstruction}}\par

已有\par

\[
M(w)=I+\dfrac{M^{(1)}}{w}+O(w^{-2}).
\]

空间 Lax 方程对去振荡矩阵 M 变成\par

\[
M_x=-\mathrm{i}\lambda\sigma_3M+A(u)M
+\mathrm{i}\rho M\sigma_3.
\]

利用 $\lambda{}=w/2+O(w^{-1})$、$\rho{}=w/2+O(w^{-1})$、$M=I+M^{(1)}/w+…$，比较 $w^{0}$ 系数：\par

\[
0=-\dfrac{\mathrm{i}}{2}\sigma_3 M^{(1)}
+A(u)+\dfrac{\mathrm{i}}{2}M^{(1)}\sigma_3.
\]

整理：\par

\[
A(u)=\dfrac{\mathrm{i}}{2}
[\sigma_3,M^{(1)}].
\]

其 (1,2) 元满足 $[\sigma{}_{3},M^{(1)}]_{12}=2M^{(1)}_{12}$，因此\par

\[
\boxed{u=\mathrm{i}M^{(1)}_{12}},
\qquad q=e^{-\mathrm{i}\Omega z}u.
\]

请记住 $M^{(1)}$ 是整个矩阵展开系数，不是 M 的第一列。第三章高能展开负责保证这个比较不是随意的形式运算。\par

\section{介质矩阵为什么要在共振点求值？}
\noindent{\small\textbf{对应原文：}\texttt{eq:RfromM}}\par

HMB 时间 Lax 对含有 $iR/(\lambda{}-a)$，所以介质变量不是只看 $w\to{}\infty{}$ 就能恢复。\par

第三章在共振点证明过：对于具有相容端态和正则基底的场，\par

\[
\chi_+^{-1}(x,a)R(x)\chi_+(x,a)
=\dfrac{\eta_0\rho(w_a)}{a}\sigma_3.
\]

两边左乘 $\chi{}_{+}$、右乘 $\chi{}_{+}^{-1}$：\par

\[
R(x)=\chi_+(x,a)
\dfrac{\eta_0\rho(w_a)}{a}\sigma_3
\chi_+(x,a)^{-1}.
\]

在规定的共振正则性条件下，把 $\chi{}$ 改写为去振荡的 M；由于中间的 $\sigma{}_{3}$ 与空间对角指数对易，指数共轭相消。于是\par

\[
\boxed{
R(x,z)=M(x,z,w_a)
\dfrac{\eta_0\rho(w_a)}{a}\sigma_3
M(x,z,w_a)^{-1}.
}
\]

为何两个 $w_{a}$ 原像给出同一 R？换叶使 $\rho{}$ 变号，同时中间 $\sigma{}_{3}$ 被交换列的 C 共轭也变号：\par

\[
C\sigma_3 C^{-1}=-\sigma_3.
\]

两个负号相消。再由 R 的分量得到\par

\[
q=e^{-\mathrm{i}\Omega z}u,\qquad
p=-e^{-\mathrm{i}\Omega z}R_{12},\qquad
\eta=R_{11}.
\]

\textbf{切记：此式在第四章首先是直接散射重构恒等式。} 若从一组抽象数据直接求得 M，还需要第五章证明 M 确实恢复完整 Lax 对，并满足共振处正则性；不能只看形式相同就断言原 PDE 成立。\par

\section{怎样直接从留数获得光场？}
\noindent{\small\textbf{对应原文：}\texttt{eq:qfinite}}\par

部分分式展开中\par

\[
\dfrac{1}{w-\zeta_j}
=\dfrac{1}{w}+O(w^{-2}).
\]

因此 $M^{(1)}$ 的 (1,2) 元来自固定 $-iQ\sigma{}_{1}/w$ 的 (1,2) 分量，以及第二组共轭留数的第一分量：\par

\[
M^{(1)}_{12}
=-\mathrm{i}Q-\sum\limits_{j=1}^{m}\overline{C_j}(u_j)_1.
\]

代入 $u=iM^{(1)}_{12}$：\par

\[
\boxed{
u=Q-\mathrm{i}\sum\limits_{j=1}^{m}\overline{C_j}(u_j)_1.
}
\]

这就是 eq:qfinite。它意味着：一旦完成有限维向量系统，光场已经能直接写出来；下一章还将同时完成介质重构和方程验证。\par

\section{本章检查清单}

\begin{enumerate}[leftmargin=2em]
\item 不看公式，从两条散射列关系重新推导 J 的四个元素。
\item 解释为什么 det M 是 \ensuremath{\gamma} 而非 1。
\item 能明确指出 M 在 w=0 的固定主项为何不能删除。
\item 对某个简单透射零点完成 Taylor 展开和留数计算。
\item 手算 h'(\ensuremath{\zeta})；说明为何配对留数不能只取复共轭。
\item 从部分分式的第二列取得 eq:finite1，再由第一列取得 eq:finite2。
\item 独立完成一次 Schur 消元。
\item 解释 theta condition 与左右端态相位的关系。
\item 用交换子计算解释 u=iM\textsuperscript{(}\textsuperscript{1}\textsuperscript{)}\textsubscript{1}\textsubscript{2}。
\item 区分“已知场的 RHP 描述”与“任意数据的逆 RHP 可解性”。
\end{enumerate}

本章数学结果的核心不是求出某个孤子的波形，而是建好了一个\textbf{统一可求解的逆问题框架}。第五章只需选择一组合法散射数据，并且补上所有必要谱约束和逆 Lax 证明。\par

